\section{Separating the scalar baseline from directional entropy}
\label{sec:critical-near-baseline-decomposition}

This section gives a structural consequence of a critical variational
path whose cost is close to the universal scalar value $V_*$. It does
not yet give a product-norm or sample-complexity bound.

\begin{lemma}[Last-exit occupation with an additive budget]
\label{lem:critical-additive-occupation}
Let $q_0,\ldots,q_H\in(0,1]$ and $D_1,\ldots,D_H\ge0$ satisfy
\[
 D_j\ge q_j-q_{j-1}+cq_{j-1}^2,\qquad
 \sum_{j=1}^iD_j\le\varepsilon+C_0q_i,
 \qquad \sum_{j=1}^H D_j\le3,
\]
where $c,C_0>0$ are fixed and $0\le\varepsilon\le1$.
Writing $m=\min_{0\le j\le H}q_j$, there is a constant $C$,
depending only on $c,C_0$, such that
\[
 H\le C\left(\frac1m+\frac{\varepsilon}{m^2}\right).
\]
Consequently, for $H\ge1$,
\[
 m\le\frac{C'}H+C'\sqrt{\frac{\varepsilon}H}.
\]
\end{lemma}

\begin{proof}
For an interval of $\ell$ transitions whose pre-states are at least
$u>0$, the dissipation inequality and concavity of $q\mapsto-1/q$
give
\[
 c\ell\le
 \frac{\text{entropy on the interval}}{u^2}
 +\frac1{q_{\rm end}}-\frac1{q_{\rm start}}.
\]
Fix a universal $u_*\in(0,1/4)$. If $m\ge u_*$, applying this
bound to the entire path gives a universal upper bound for $H$,
which suffices. Otherwise choose the least $J$ with
$u_J=2^Jm\ge u_*$. Define the nondecreasing last-exit indices
\[
 \tau_j=\max\{i:q_i\le u_j\},\qquad 0\le j\le J.
\]
The prefix ending at $\tau_0$ has entropy at most
$\varepsilon+C_0m$ and all its states are at least $m$, so its
length is at most $C(\varepsilon/m^2+1/m)$.

For a nonempty block from $\tau_j$ to $\tau_{j+1}$, remove its
first transition. All remaining pre-states exceed $u_j$, and the
endpoint lies in $(u_j,2u_j]$. The entropy of this interval is at
most the full prefix entropy at $\tau_{j+1}$, hence at most
$\varepsilon+2C_0u_j$. Its length is therefore at most
$C(\varepsilon/u_j^2+1/u_j)$. Restoring the removed step adds one.
After $\tau_J$, remove one first step if present; the remaining
pre-states exceed $u_J\ge u_*$ and the remaining entropy is at
most three, so this tail has universally bounded length.
Summing the geometric series in $u_j^{-1}$ and $u_j^{-2}$,
and absorbing the at most $J$ removed transitions, proves the first
claim. Solving $Hm^2\le C(m+\varepsilon)$ proves the second.
\end{proof}

\begin{theorem}[Near-baseline directional budget]
\label{thm:critical-directional-budget}
Consider an admissible critical variational path of depth $L\ge2$
with hidden second moments $Q_0,\ldots,Q_H$, $H=L-1$, hidden
entropy costs $D_1,\ldots,D_H$, output $z$, and readout cost
\[
 r=\tfrac12z^{\mathsf T}Q_H^+z.
\]
Assume its total cost $A=\sum_{j=1}^H D_j+r$ satisfies
$A\le V_*+\varepsilon$ with $0\le\varepsilon\le1$ and $A\le3$,
and that $f(x)z(x)\ge1/2$ on the Boolean cube.

Choose a bichromatic cube edge if $f$ is nonconstant, and an
antipodal pair if $f$ is constant. In the first case set
$b=(e_x-e_y)/2$; in the second set $b=(e_x+e_y)/2$.
Write
\[
 v_j=b^{\mathsf T}Q_jb,\qquad
 a=b^{\mathsf T}z,\qquad
 m=\min_{0\le j\le H}v_j,
\]
and choose an index $i$ attaining the minimum. All $v_j$ are
positive, $v_0=1/d$, and
\[
 m\le\frac C H+C\sqrt{\frac{\varepsilon}H}.
\]

For each $j>i$, lift the controlled feature law through the inverse
coordinatewise $\tanh$ map to a preactivation law $\rho_j$.
Let $\eta_j$ be its pushforward under $G\mapsto b^{\mathsf T}G$,
and define the conditional entropy
\[
 E_j=D(\rho_j\Vert N(0,Q_{j-1}))
       -D(\eta_j\Vert N(0,v_{j-1}))\ge0.
\]
Then the nonnegative quantities
\[
 r_\perp=r-\frac{a^2}{2v_H},\qquad
 r_{\rm amp}=\frac{a^2-1/4}{2v_H}
\]
satisfy
\[
 \boxed{\quad
 \sum_{j=1}^iD_j+\sum_{j=i+1}^H E_j
       +r_\perp+r_{\rm amp}
 \le\varepsilon+C_0m.
 \quad}
\]
In particular this budget is at most
$C\{\varepsilon+H^{-1}+\sqrt{\varepsilon/H}\}$.
More precisely, define the nonnegative scalar Bellman slacks
\[
 s_j=D(\eta_j\Vert N(0,v_{j-1}))+V(v_j)-V(v_{j-1}),\qquad
 s_{\rm out}=\frac1{8v_H}-V(v_H).
\]
They satisfy the exact decomposition
\[
 A-V(m)=\sum_{j=1}^iD_j+\sum_{j=i+1}^H(E_j+s_j)
          +r_\perp+r_{\rm amp}+s_{\rm out}.
\]
Thus all the displayed nonnegative terms together obey the same
upper bound $\varepsilon+C_0m$.
\end{theorem}

\begin{proof}
The chosen pair has $|a|\ge1/2$ and input variance $v_0=1/d$.
The signed-pair comparison in
Lemma~\ref{lem:critical-pair-scalar-comparison} gives, for every
prefix,
\[
 \sum_{j=1}^tD_j\le A-V(v_t)
 \le\varepsilon+C_0v_t.
\]
Here $V_*-V(v)\le C_0v$ extends from small $v$ to $(0,1]$
by increasing $C_0$, since $0\le V(v)\le V_*<2$.
The same comparison gives the dissipation inequality, and the full
hidden entropy is at most $A\le3$. Thus
Lemma~\ref{lem:critical-additive-occupation} gives the claimed
minimum-variance bound. Positivity follows because a zero signed-pair
variance propagates under every finite-entropy transition and is
incompatible with $|a|\ge1/2$.

Coordinatewise $\tanh$ is a measurable bijection onto its image, so
lifting preserves relative entropy. The chain rule under the scalar
projection gives $D_j=D_j^{\rm sc}+E_j$, where
$D_j^{\rm sc}=D(\eta_j\Vert N(0,v_{j-1}))$; equivalently, $E_j$
is the average KL divergence of the controlled conditional law given
$b^{\mathsf T}G$ from the corresponding conditional Gaussian law.
This identity remains valid for singular covariances by disintegration
on their Gaussian support.

The signed-pair contraction inequality gives
$\int\tanh^2s\,d\eta_j(s)\ge v_j$. The scalar Bellman inequality
and monotonicity of $V$ therefore imply
\[
 D_j^{\rm sc}\ge V(v_{j-1})-V(v_j).
\]
This proves $s_j\ge0$, while the length-zero scalar path gives
$s_{\rm out}\ge0$.
After summing over $j>i$ and adding $1/(8v_H)\ge V(v_H)$,
\[
 \sum_{j=i+1}^H D_j^{\rm sc}+\frac1{8v_H}\ge V(v_i).
\]
The finite readout cost implies $z\in\operatorname{range}Q_H$.
Cauchy--Schwarz on that range gives $r_\perp\ge0$, and
$r_{\rm amp}\ge0$ follows from $|a|\ge1/2$. Since
\[
 r=\frac1{8v_H}+r_\perp+r_{\rm amp},
\]
the total path cost is at least
\[
 \sum_{j=1}^iD_j+\sum_{j=i+1}^H E_j
       +r_\perp+r_{\rm amp}+V(v_i).
\]
Subtract $V(v_i)$ and use $A\le V_*+\varepsilon$ and $v_i=m$
to obtain the boxed inequality. Retaining the scalar slacks in the
same telescoping calculation gives the asserted exact decomposition.
\end{proof}

\begin{corollary}[Exact classification by one anchored kernel section]
\label{cor:critical-anchored-section}
Under the theorem's hypotheses set $B=\varepsilon+C_0m$.
If $B<1/8$, then on the entire Boolean cube
\[
 f=\operatorname{sign}\!\left(\frac a{v_H}Q_Hb\right).
\]
If $B\le1/32$, then the signed section
$\operatorname{sign}(a)Q_Hb$ has uniform margin at least
$1/(16A)\ge1/48$ for $f$.
\end{corollary}

\begin{proof}
Put $z_\perp=z-(a/v_H)Q_Hb$. Orthogonal decomposition in the
$Q_H^+$ seminorm gives
\[
 \tfrac12z_\perp^{\mathsf T}Q_H^+z_\perp=r_\perp\le B.
\]
Since $Q_H(u,u)\le1$, coordinatewise Cauchy--Schwarz gives
$|z_\perp(u)|\le\sqrt{2B}$. This is strictly less than the output
margin $1/2$ when $B<1/8$, proving the sign identity. If $B\le1/32$,
the anchored output has margin at least $1/4$. Finally
$a^2/(2v_H)\le A$ and $|a|\ge1/2$ imply
$|a|/v_H\le4A$. Dividing the anchored output by that factor gives
the asserted margin for the unscaled signed section.
\end{proof}

\begin{proposition}[Zero directional entropy gives a shallow classifier]
\label{prop:critical-zero-directional-entropy}
Under the theorem's hypotheses, suppose additionally that
\[
 \sum_{j=1}^iD_j=0,\qquad
 \sum_{j=i+1}^H E_j=0,\qquad r_\perp<1/8.
\]
Then $f$ has a depth-two $\tanh$ implementation with unit global
margin, width polynomial in $d$, and product of layer Frobenius
norms polynomial in $d$. The bounds are independent of $L$.
\end{proposition}

\begin{proof}
Let $\mathcal G$ be the group of signed coordinate permutations
fixing the two selected cube inputs individually. It fixes the vector
$b$ under its induced permutation of input indices. The raw input
covariance is invariant under this group, and the zero-entropy prefix
is exactly the uncontrolled feature recursion, so $Q_i$ is invariant.

At a later layer, $E_j=0$ means that the controlled preactivation
law is obtained by changing only the distribution of
$S=b^{\mathsf T}G$ and retaining the reference conditional Gaussian
law given $S$. If the reference covariance is $\mathcal G$ invariant,
so is this conditional construction, because $S$ is fixed by every
group element. Coordinatewise $\tanh$ preserves the invariance.
Induction therefore makes $Q_H$ invariant. The anchored section
$Q_Hb$ is consequently constant on every $\mathcal G$ orbit.
The readout residual bound, by the preceding proof, identifies its
sign with $f$.

After an input sign change, an edge pair is
$(1,\ldots,1)$ and $(-1,1,\ldots,1)$. Permuting coordinates
$2,\ldots,d$ shows that $f$ depends only on $x_1$ and
$\sum_{j=2}^d x_j$, giving at most $2d$ orbit classes.
These classes have distinct, separated integer codes
\[
 t(x)=(d+1)x_1+\sum_{j=2}^d x_j.
\]
For an antipodal pair the corresponding orbit classes are determined
by $t(x)=\sum_{j=1}^d x_j$ and number at most $d+1$.
In either case the codes have magnitude $O(d)$ and gaps at least one.

Order their distinct values $t_1<\cdots<t_M$, and let $f_j$ be the
desired sign at $t_j$. With thresholds halfway between consecutive
values, the exact step expansion is
\[
 f(x)=c_0+\sum_{j=1}^{M-1}c_j
             \operatorname{sign}(t(x)-\theta_j),
 \qquad c_j=(f_{j+1}-f_j)/2,\quad c_0=(f_1+f_M)/2.
\]
All these coefficients have absolute value at most one.
Replace each step by $\tanh(A(t-\theta_j))$ with
$A=C\log(d+2)$. The sum's uniform error is at most $1/4$ for a
sufficiently large universal $C$. A constant first-layer feature
implements $c_0$ exactly by adjusting its readout coefficient.
The resulting width, incoming coefficients, and readout norm are
polynomial in $d$, and doubling the readout gives unit margin.
Input sign changes preserve Frobenius norms. This proves the claim.
\end{proof}

The zero-entropy proposition is a rigidity statement. A quantitative
version controlling intrinsic product norm by
$\operatorname{poly}(d,L)\exp\{C\sqrt{L\varepsilon}\}$ would require
stability of the conditional-Gaussian suffix under the small
directional budget. That implication is not proved here.
